\chapter[Прием QPSK BPSK, прием на согласованный фильтр, глазковая диаграмма, поиск оптимального отсчетного значения]{Прием QPSK BPSK, прием на согласованный фильтр, глазковая диаграмма, поиск оптимального отсчетного значения и необходимость символьной синхронизации}

\begin{leftbar}
    \href{https://github.com/someotb/sdr}{Ссылка на GitHub}
\end{leftbar}

\section*{Цель практики:}

Сформировать биты, пропустить их через согласованный фильтр. Построить глазковую диаграмму, в ручную найти оптимальное отсчетное значение
 
\section*{Краткие теоретические сведения}

\textbf{Согласованный фильтр (matched filter)} в цифровой обработке сигналов — это линейный фильтр, обеспечивающий максимальное отношение сигнал/шум (SNR) на выходе в момент принятия решения при наличии аддитивного белого гауссовского шума (AWGN).  

Для известного сигнала \( s(t) \) длительностью \( T \), импульсная характеристика согласованного фильтра задаётся как:
$$
h(t) = s(T - t),
$$
то есть представляет собой зеркально отражённую во времени копию сигнала \( s(t) \), сдвинутую на \( T \).

В дискретной форме, для последовательности отсчётов \( s[n] \), \( n = 0, 1, \dots, N-1 \), импульсная характеристика согласованного фильтра:
$$
h[n] = s^*[N - 1 - n],
$$
где \( (\cdot)^* \) обозначает комплексное сопряжение.

Выход согласованного фильтра вычисляется как свёртка входного сигнала \( x[n] \) с \( h[n] \):
$$
y[n] = \sum_{k=0}^{N-1} x[k] \cdot s^*[k - n],
$$
что эквивалентно корреляции входного сигнала с ожидаемой формой сигнала.  

Согласованный фильтр широко применяется в программно-определяемых радиосистемах (SDR) для обнаружения и демодуляции сигналов с известной структурой (например, пилотных символов, преамбулы или кодов расширения в CDMA).

\vspace{20pt}

\textbf{Глазковая диаграмма (eye diagram)} — это визуальное представление цифрового сигнала, полученное путём наложения друг на друга нескольких последовательных тактовых интервалов в осциллографе или программной среде.  

Формируется глазковая диаграмма следующим образом: осциллограф или программа захватывает отрезки сигнала длительностью, кратной периоду тактовой частоты \( T \), и отображает их синхронно по фронту тактового сигнала. В результате наложения траекторий сигнала возникает характерная форма, напоминающая «глаз».

Качество цифровой передачи оценивается по степени раскрытия «глаза»:
- чем шире и выше глаз, тем меньше межсимвольная интерференция (ISI) и шум;
- закрытый глаз указывает на сильную ISI, джиттер или высокий уровень шума.

Глазковая диаграмма (eye diagramm) дает наглядное представление о форме
принимаемого сигнала на выходе согласованного фильтра. Глазковая диаграмма строится в виде накапливающего изображения сигналов на выходе СФ в
интервале времени

$$
0.5T_s < t < 0.5T_S
$$

по большому числу символов

Если изобразить много таких отрезков сигнала на выходе СФ в виде накапливающейся суммы изображений, то получим глазковую диаграмму.

\begin{figure}[H]
    \centering
    \includegraphics[width=1.0\textwidth]{eye_teory.png}
    \caption{Глазковая диаграмма}
\end{figure}

\section*{Выполнение}

Согласованный фильтр уже реализован в функции filter

\begin{lstlisting}[language=C++, caption=filter]
void filter(complex<double> *symbols_ups, int len_symbols_ups, complex<double> *impulse, int L) {
    vector<complex<double>> sum(len_symbols_ups, 0);
    for (size_t i = 0; i < (size_t)len_symbols_ups; i++) {
        for (size_t j = 0; j < (size_t)L && (int)(i-j) >= 0; j++) {
            sum[i] += impulse[j] * symbols_ups[i-j];
        }
    }
    memcpy(symbols_ups, sum.data(), len_symbols_ups * sizeof(complex<double>));
}
\end{lstlisting}

Чтобы получить глазковую диаграмму, нужно написать наш код на python

\begin{lstlisting}[language=python, caption=python eye diagram]
import numpy as np
import matplotlib.pyplot as plt

data = np.fromfile("/home/someotb/Files/Code/sdr/pluto/dev/build/rx.pcm", dtype=np.int16)

samples = []
for x in range(0, len(data) - 1, 2):
    samples.append(data[x] + 1j * data[x+1])

real = np.real(samples)

sps = 10
segments = []
for i in range(0, len(real) - 2*sps, sps):
    segments.append(real[i:i + 2*sps])

plt.plot(np.array(segments).T)
plt.xlim(7.5, 12.5)
plt.grid(True)
plt.show()
\end{lstlisting}

\vspace{20pt}

Получим такой результат

\vspace{20pt}

\begin{figure}[H]
    \centering
    \includegraphics[width=0.8\textwidth]{eyediagram.png}
    \caption{Глазковая диаграмма}
\end{figure}

Поиск оптимального значения отсчетного значения мы будем выполнять вручную. Сперва напишем код для отображения сигнального созвездия

\begin{lstlisting}[language=python, caption=python eye diagram]
from matplotlib import pyplot as plt
import numpy as np

rx = np.fromfile(f"/home/someotb/Files/Code/sdr/pluto/dev/build/rx.pcm", dtype=np.int16)

samples = []

for x in range(0, len(rx), 2):
    samples.append((rx[x] + 1j * rx[x+1])/np.max(rx))
    
impulse = np.ones(10)

I = np.real(samples)
Q = np.imag(samples)

i_convolve = np.convolve(I, impulse)
q_convolve = np.convolve(Q, impulse)
i_con_dec = []
q_con_dec = []

for i in range(-1, len(i_convolve), 10):
    i_con_dec.append(i_convolve[i])
    q_con_dec.append(q_convolve[i])

# Убираем шумы
i_con_dec = [x for x in i_con_dec if not (abs(x) <= 1)]
q_con_dec = [x for x in q_con_dec if not (abs(x) <= 1)]
    
plt.scatter(i_con_dec, q_con_dec)
plt.axhline()
plt.axvline()
plt.show()
\end{lstlisting}

Строчки когда я убираю шумы, иногда меняют длину массивов, что приводит к ошибке, поэтому иногдя я их комментирую, и для таких значений
у сигнального созвездия будет шум в центре

\begin{figure}[H]
    \centering
    \includegraphics[width=1.0\textwidth]{_1.png}
    \caption{Отсчетное значение = -1}
\end{figure}

\begin{figure}[H]
    \centering
    \includegraphics[width=1.0\textwidth]{1.png}
    \caption{Отсчетное значение = 1}
\end{figure}

\begin{figure}[H]
    \centering
    \includegraphics[width=1.0\textwidth]{5.png}
    \caption{Отсчетное значение = 5}
\end{figure}


\begin{figure}[H]
    \centering
    \includegraphics[width=1.0\textwidth]{7.png}
    \caption{Отсчетное значение = 7}
\end{figure}

По итогу самое лучшее отсчетное значение оказалось -1, сигнальное созвездие с этим значением выглядит лучше всего, более похожее на сигнальное созвездие
QPSK

\section*{Вывод}
Получили глазковую диаграмму, в ручную подобрали отсчетное значения для лучшей схожести с сигнальным созвездием QPSK, получили графики для
разных отсчетных значений.